% ============================================================================
%  Chapter 10 -- Introduction to Electronics and Digital Electronics (Lecture 10)
% ============================================================================
\chapter{Electronics: What a Bit Is Made Of}
\label{ch:electronics}

\epigraphline{Any sufficiently advanced technology is indistinguishable from
magic.}{Arthur C. Clarke, \emph{Profiles of the Future}, 1973}

\begin{chapgoals}
By the end of this chapter you will be able to
\begin{itemize}
\item use voltage, current, resistance and power, and apply Ohm's law to a simple circuit;
\item size the resistor of an LED and compute a voltage divider;
\item explain what a transistor is and how a MOS transistor works as a switch;
\item say how voltages become logic values, and what a noise margin is;
\item draw the CMOS inverter, NAND and NOR, and count their transistors;
\item explain pull-up resistors, floating inputs and why a gate takes time and energy to switch.
\end{itemize}
\end{chapgoals}

The gates of the last two chapters were symbols on paper. This chapter gives
them a physical body, before the next one composes them into machines. No
previous knowledge of physics is assumed. The numbers quoted for the kit come
from the datasheet of the ESP32-C6 \citep{esp32c6trm} and from the hardware
documentation of \prg.

% ----------------------------------------------------------------------------
\section{Warm-up on the machine}

\asmfile[caption={One button, one bit.},label=lst:el:button]{code/l10/s04.S}

Since the first chapters a pressed button has been a 1 in a register
(Listing~\ref{lst:el:button}). On the board the button is a piece of metal
that joins a pin of the chip to ground. Pressed means 0~volts on the pin, and
yet the program reads 1. What is a volt? Why is the released pin at 3.3~V?
And who decides that 0~V is a bit? The chapter follows the button from the
finger to the bit, and its last section answers the three questions.

% ----------------------------------------------------------------------------
\section{Electricity in four quantities}

\subsection{Voltage, current, resistance, power}

Electric charge is carried by electrons, and four quantities describe what it
does in a circuit (Table~\ref{tab:el:four}). A comparison with water in pipes
helps the intuition.

\begin{table}[htb]
\centering\small
\caption{Four quantities.}\label{tab:el:four}
\begin{tabularx}{\textwidth}{@{}llLL@{}}
\toprule
\textbf{Quantity} & \textbf{Symbol, unit} & \textbf{Water analogy} & \textbf{On the kit}\\
\midrule
voltage & $V$, volt (V) & the pressure that pushes the water & 3.3 V between the supply rail and ground\\
current & $I$, ampere (A) & the flow: litres per second & about 10 mA through a lit LED\\
resistance & $R$, ohm ($\Omega$) & a narrow pipe that hinders the flow & 45 k$\Omega$: the pull-up inside each input pin\\
power & $P$, watt (W) & the work the flow can do each second & at most 2.5 W from a USB 2.0 port (5 V, 0.5 A)\\
\bottomrule
\end{tabularx}
\end{table}

A voltage is always measured \emph{between two points}. When one says that
``the pin is at 3.3~V'' one means with respect to \term{ground} (GND), the
common reference of the whole circuit, which is at 0~V by definition. The
usual prefixes are milli (m, $10^{-3}$), micro ($\mu$, $10^{-6}$), nano (n,
$10^{-9}$), pico (p, $10^{-12}$), kilo (k, $10^{3}$) and mega (M, $10^{6}$).

\subsection{Ohm's law, and the resistor of an LED}

\begin{theorem}{Ohm's law, 1827}
The current through a resistor is proportional to the voltage across it:
\[ V = R\cdot I . \]
\end{theorem}

A circuit is a closed loop: the current leaves the source, crosses the
components and returns. An LED (Figure~\ref{fig:el:led}) lights when current
flows through it, and when it does there is a drop of about 2~V across it;
without a resistor in series the current would be far too large. The resistor
takes what is left, $3.3-2=1.3$~V. For a current of 10~mA,
\[ R=\frac{1.3\ \mathrm{V}}{0.010\ \mathrm{A}}=130\ \Omega, \]
and the nearest common value is 150~$\Omega$, which gives
$1.3/150 = 8.7$~mA. The forward voltage of an LED depends on its colour:
about 1.8--2.2~V for red, about 3~V for blue and white.

\begin{figure}[htb]
\centering
\begin{subfigure}[b]{0.3\textwidth}\centering
\begin{circuitikz}[scale=0.85,transform shape]
  \draw (0,0) node[ground]{} to[battery1,l=3.3<\volt>,invert] (0,3) -- (2.2,3)
    to[R,l=150<\ohm>] (2.2,1.6) to[leDo,l=LED] (2.2,0) -- (0,0);
\end{circuitikz}
\caption{An LED and its resistor}\label{fig:el:led}
\end{subfigure}\hfill
\begin{subfigure}[b]{0.3\textwidth}\centering
\begin{circuitikz}[scale=0.85,transform shape]
  \draw (0,3.2) node[vcc]{$V_{dd}$} to[R,l=$R_1$] (0,1.6) to[R,l=$R_2$] (0,0) node[ground]{};
  \draw (0,1.6) node[circ]{} -- (1.2,1.6) node[right]{$V_{out}$};
\end{circuitikz}
\caption{A voltage divider}\label{fig:el:div}
\end{subfigure}\hfill
\begin{subfigure}[b]{0.34\textwidth}\centering
\begin{circuitikz}[scale=0.85,transform shape]
  \draw (0,0.2) node[nmos](n){} (n.G) node[left]{G} (n.D) node[above]{D} (n.S) node[below]{S};
  \draw (2.4,0.2) node[pmos](p){} (p.G) node[left]{G} (p.S) node[above]{S} (p.D) node[below]{D};
  \node[lbl] at (0,-1.3) {nMOS};
  \node[lbl] at (2.4,-1.3) {pMOS};
\end{circuitikz}
\caption{The two MOS transistors}\label{fig:el:mos}
\end{subfigure}
\caption{Three elementary circuits.}
\end{figure}

\subsection{Series, parallel and the voltage divider}

Two resistors are in \term{series} when the same current crosses both: the
resistances add, $R=R_1+R_2$. They are in \term{parallel} when the same
voltage is across both: the current splits, and the total resistance,
$1/R=1/R_1+1/R_2$, is smaller than either.

Two resistors in series divide the voltage (Figure~\ref{fig:el:div}):
\begin{equation}\label{eq:el:div}
V_{out} = V_{dd}\,\frac{R_2}{R_1+R_2}.
\end{equation}
Equal resistors give half the supply; a much smaller $R_2$ pulls $V_{out}$
close to ground. The formula follows from Ohm's law, since the current
$I=V_{dd}/(R_1+R_2)$ crosses $R_2$, and from Kirchhoff's rules: the currents
that enter a node add up to those that leave it, and the voltages round a loop
add up to zero. The divider is the key to the circuit of the button, where
$R_2$ is a switch: a resistance that is either infinite or nearly zero.

\subsection{Power and energy}

\term{Power} is voltage times current, $P=V\cdot I$, in watts; in a resistor
it becomes heat, $P=R\cdot I^{2}$. \term{Energy} is power times time, in
joules; a battery stores energy. The LED above consumes
$3.3\ \mathrm{V}\times 8.7\ \mathrm{mA} = 29$~mW, of which 11~mW heat the
resistor. Every transistor that switches inside a processor dissipates a
little energy, and with thousands of millions of them heat is the limit of
modern chips: a laptop processor dissipates tens of watts in a few square
centimetres.

\subsection{The capacitor: why signals take time}

A \term{capacitor} stores electric charge as a tank stores water; its
capacitance is measured in farads (F). Its voltage cannot jump: the capacitor
must be filled, or emptied, through a resistance, and the time that this
takes is of the order of the product
\[ \tau = R\cdot C, \]
the \term{time constant}, after which the voltage has covered about 63\% of
the way. An input pin of the kit has a capacitance of about 2~pF and rises
through its pull-up of 45~k$\Omega$ in about
$45\ \mathrm{k}\Omega\times 2\ \mathrm{pF}=90$~ns (the wiring and the button
add more, so the real figure is larger).

Every wire and every input of a transistor is a small capacitor. This is
where the propagation delay of a gate comes from: resistance times
capacitance is a time, and it is the reason why nothing in a circuit is
instantaneous. A capacitor that is charged or empty is also a way of storing
a bit: the DRAM cell of Chapter~\ref{ch:memory}.

% ----------------------------------------------------------------------------
\section{Semiconductors}

\subsection{Conductors, insulators, semiconductors}

Conductors, such as copper, let charge flow freely; insulators, such as
glass, do not. \term{Semiconductors}, above all silicon, lie in between, and
their conduction can be controlled. The control is obtained by
\term{doping}: adding a few atoms of another element. Silicon of type~n has
extra electrons, which are negative carriers of charge; silicon of type~p has
missing electrons, called \emph{holes}, which behave as positive carriers.
Where a p region meets an n region there is a \term{junction}, the building
block of every device. A chip is a slice of silicon with thousands of
millions of doped regions, connected by layers of metal. Silicon is the
second most abundant element in the Earth's crust, the element of sand, and a
whole industry rests on purifying it and doping it with extreme precision.

\subsection{The diode and the LED}

A single junction makes a \term{diode}, a one-way valve for current. In the
forward direction, from p to n, it conducts once the voltage exceeds a
threshold, about 0.7~V for silicon; in the reverse direction it blocks. A
light-emitting diode gives off light when it conducts; its colour depends on
the material, and so does its threshold. Diodes turn alternating current into
direct current in every power supply. LEDs are everywhere in a computer:
indicators, the backlight of the display of Chapter~\ref{ch:video}, optical
links; and the laser of the optical disks of Chapter~\ref{ch:storage} is a
diode too.

\subsection{The transistor: a switch with no moving parts}

A \term{transistor} is a device with three terminals in which a small signal
on one controls the current between the other two. Used as an amplifier it
made the radio portable; used as a switch it made the computer
(Table~\ref{tab:el:history}). Vacuum tubes and relays did the same job in the
first computers of Chapter~\ref{ch:history}, a million times larger and
slower.

\begin{table}[htb]
\centering\small
\caption{The transistor in six dates.}\label{tab:el:history}
\begin{tabularx}{\textwidth}{@{}lP{0.34\textwidth}L@{}}
\toprule
\textbf{Year} & \textbf{Who} & \textbf{What}\\
\midrule
1947 & Bardeen, Brattain and Shockley, Bell Labs & the first transistor: it replaces the vacuum tube\\
1958--59 & Jack Kilby; Robert Noyce & the integrated circuit: several transistors on one piece of semiconductor\\
1959--60 & Atalla and Kahng, Bell Labs & the MOS transistor: small, simple to build, almost no control current\\
1963 & Wanlass and Sah, Fairchild & CMOS: pairs of complementary MOS transistors that draw power only when switching\\
1971 & Intel 4004 & 2\,300 MOS transistors: a processor on a chip\\
today & the whole industry & tens of thousands of millions of CMOS transistors on a chip\\
\bottomrule
\end{tabularx}
\end{table}

\subsection{The MOS transistor as a switch}

The name MOS lists the three layers of the device: metal, oxide,
semiconductor. Its three terminals are the \term{gate}, which controls
whether current can pass, and the \term{source} and the \term{drain}, between
which it passes (Figure~\ref{fig:el:mos}). The gate is insulated: it draws
practically no current (at most 50~nA on a pin of the ESP32-C6), only the tiny
charge of its capacitance.

\begin{definition}{nMOS and pMOS}
An \term{nMOS} transistor conducts (the switch is closed) when its gate is
high, and blocks when it is low. A \term{pMOS} transistor conducts when its
gate is low, and blocks when it is high.
\end{definition}

The two types are complementary: what closes one opens the other. The small
circle on the gate of the pMOS means, as on logic symbols, ``active when
low''. One more fact decides how they are used: an nMOS passes a good 0, and
is therefore used to connect to ground; a pMOS passes a good 1, and is used
to connect to the supply.

% ----------------------------------------------------------------------------
\section{From volts to bits}

\subsection{Logic levels: when is a voltage a 1?}

A digital circuit uses two \emph{ranges} of voltage: low is 0, high is 1
(Figure~\ref{fig:el:levels}). An output guarantees clean levels: on the kit,
above 2.64~V for a 1 and below 0.33~V for a 0. An input accepts wider ranges:
above 2.475~V it reads 1, below 0.825~V it reads 0. (With $V_{DD}=3.3$~V the
datasheet gives $V_{OH}\ge 0.8\,V_{DD}$, $V_{OL}\le 0.1\,V_{DD}$,
$V_{IH}\ge 0.75\,V_{DD}$, $V_{IL}\le 0.25\,V_{DD}$.)

\begin{definition}{Noise margin}
The \term{noise margin} is the difference between the level that an output
guarantees and the level that an input requires. A disturbance smaller than
the noise margin cannot change the bit.
\end{definition}

On the kit the margins are $2.64-2.475=0.165$~V for the high level and
$0.825-0.33=0.495$~V for the low level. Between the two input thresholds the
value is undefined: signals must cross that band quickly and never stay there.

\begin{figure}[htb]
\centering
\begin{tikzpicture}[x=1cm,y=1.5cm]
  % output column
  \fill[prgaccent!30] (0,2.64) rectangle (2.2,3.3);
  \fill[prgaccent!30] (0,0) rectangle (2.2,0.33);
  \draw[prgsteel] (0,0) rectangle (2.2,3.3);
  \node[lbl] at (1.1,3.55) {\textbf{output}};
  \node[lbl] at (1.1,2.97) {1};
  \node[lbl] at (1.1,0.16) {0};
  % input column
  \fill[prgamber!30] (5,2.475) rectangle (7.2,3.3);
  \fill[prgamber!30] (5,0) rectangle (7.2,0.825);
  \fill[prgred!12] (5,0.825) rectangle (7.2,2.475);
  \draw[prgsteel] (5,0) rectangle (7.2,3.3);
  \node[lbl] at (6.1,3.55) {\textbf{input}};
  \node[lbl] at (6.1,2.9) {reads 1};
  \node[lbl] at (6.1,0.41) {reads 0};
  \node[lbl,text=prgred] at (6.1,1.65) {undefined};
  \foreach \y/\l in {3.3/{3.3 V},2.64/{2.64 V},0.33/{0.33 V},0/{0 V}} \node[lbl,anchor=east] at (-0.1,\y) {\l};
  \foreach \y/\l in {2.475/{2.475 V},0.825/{0.825 V}} \node[lbl,anchor=west] at (7.3,\y) {\l};
  \draw[dashed,prgink] (2.2,2.64) -- (5,2.64); \draw[dashed,prgink] (2.2,2.475) -- (5,2.475);
  \draw[dashed,prgink] (2.2,0.33) -- (5,0.33); \draw[dashed,prgink] (2.2,0.825) -- (5,0.825);
  \draw[<->,prgviolet,thick] (3.6,2.475) -- (3.6,2.64) node[midway,right,lbl,text=prgviolet]{margin 0.165 V};
  \draw[<->,prgviolet,thick] (3.6,0.33) -- (3.6,0.825) node[midway,right,lbl,text=prgviolet]{margin 0.495 V};
\end{tikzpicture}
\caption{Logic levels of the ESP32-C6 with a supply of 3.3 V.}
\label{fig:el:levels}
\end{figure}

\subsection{Why digital?}

In an \emph{analog} circuit the exact value of the voltage is the
information. Every stage adds a little noise and distortion, which
accumulate, and copies get worse at every generation. Analog circuits remain
indispensable at the edges of a system: sensors, radio, sound. In a
\emph{digital} circuit only the range matters, high or low. Every gate
restores the levels, so that noise does not accumulate; copies are perfect,
and errors can be detected and corrected (Chapter~\ref{ch:combinational}).
And two levels are easy to tell apart with transistors that are fast, tiny
and cheap. It is the story of Chapter~\ref{ch:information}, seen from the
wire.

\subsection{The CMOS inverter: a NOT gate from two transistors}

The simplest gate is made of one pMOS towards the supply and one nMOS towards
ground, driven by the same input (Figure~\ref{fig:el:inv}).

\begin{figure}[htb]
\centering
\begin{subfigure}[b]{0.28\textwidth}\centering
\begin{circuitikz}[scale=0.8,transform shape]
  \draw (0,2.0) node[pmos](p){} (0,0) node[nmos](n){};
  \draw (p.S) -- ++(0,0.3) node[vcc]{3.3 V};
  \draw (n.S) node[ground]{};
  \draw (p.D) -- (n.D);
  \draw (p.G) -- ++(-0.4,0) coordinate (g1) |- (n.G);
  \draw (g1) ++(0,-1.0) node[circ]{} -- ++(-0.6,0) node[left]{$A$};
  \draw (0,1.0) node[circ]{} -- ++(0.9,0) node[right]{$Y$};
\end{circuitikz}
\caption{NOT: $Y=\nott{A}$}\label{fig:el:inv}
\end{subfigure}\hfill
\begin{subfigure}[b]{0.34\textwidth}\centering
\begin{circuitikz}[scale=0.8,transform shape]
  \draw (0,3.6) node[pmos](p1){} (2,3.6) node[pmos](p2){};
  \draw (1,1.6) node[nmos](n1){} (1,0) node[nmos](n2){};
  \draw (p1.S) -- ++(0,0.25) coordinate (v1) (p2.S) -- ++(0,0.25) coordinate (v2);
  \draw (v1) -- (v2); \draw (1,0|-v1) node[vcc]{3.3 V};
  \draw (p1.D) -- ++(0,-0.25) coordinate (o1) (p2.D) -- ++(0,-0.25) coordinate (o2);
  \draw (o1) -- (o2);
  \draw (n1.D) -- (n1.D|-o1) node[circ]{};
  \draw (o2) -- ++(0.7,0) node[right]{$Y$};
  \draw (n1.S) -- (n2.D); \draw (n2.S) node[ground]{};
  \node[left] at (p1.G) {$A$}; \node[left] at (p2.G) {$B$};
  \node[left] at (n1.G) {$A$}; \node[left] at (n2.G) {$B$};
\end{circuitikz}
\caption{NAND: $Y=\nott{A\cdot B}$}\label{fig:el:nand}
\end{subfigure}\hfill
\begin{subfigure}[b]{0.34\textwidth}\centering
\begin{circuitikz}[scale=0.8,transform shape]
  \draw (1,5.2) node[pmos](p1){} (1,3.6) node[pmos](p2){};
  \draw (0,1.4) node[nmos](n1){} (2,1.4) node[nmos](n2){};
  \draw (p1.S) -- ++(0,0.25) node[vcc]{3.3 V};
  \draw (p1.D) -- (p2.S);
  \draw (n1.D) -- ++(0,0.25) coordinate (o1) (n2.D) -- ++(0,0.25) coordinate (o2);
  \draw (o1) -- (o2);
  \draw (p2.D) -- (p2.D|-o1) node[circ]{};
  \draw (o2) -- ++(0.7,0) node[right]{$Y$};
  \draw (n1.S) -- ++(0,-0.2) coordinate (g1) (n2.S) -- ++(0,-0.2) coordinate (g2);
  \draw (g1) -- (g2); \draw (1,0|-g1) node[ground]{};
  \node[left] at (p1.G) {$A$}; \node[left] at (p2.G) {$B$};
  \node[left] at (n1.G) {$A$}; \node[left] at (n2.G) {$B$};
\end{circuitikz}
\caption{NOR: $Y=\nott{A+B}$}\label{fig:el:nor}
\end{subfigure}
\caption{Three CMOS gates. In the NAND and in the NOR the gates of the
transistors marked with the same letter are connected together.}
\label{fig:el:cmos}
\end{figure}

\begin{center}\small
\begin{tabular}{@{}cccl@{}}
\toprule
$A$ & \textbf{pMOS} & \textbf{nMOS} & $Y$\\
\midrule
0 & on & off & 1: joined to 3.3 V\\
1 & off & on & 0: joined to ground\\
\bottomrule
\end{tabular}
\end{center}

Exactly one of the two transistors conducts: the output is always connected to
one rail, never to both. With the input still, no current flows from the
supply to ground, so that CMOS (\emph{complementary} MOS) consumes almost
nothing at rest. And the output is a fresh, clean level whatever the quality
of the input: the gate \emph{regenerates} the signal.

\subsection{NAND and NOR in CMOS}

\begin{prgconcept}[Series is AND, parallel is OR]
Two switches in series conduct when the first \emph{and} the second are
closed; two switches in parallel conduct when the first \emph{or} the second
is closed. It is the idea of Nakajima and Shannon.
\end{prgconcept}

In the NAND (Figure~\ref{fig:el:nand}) two nMOS in series reach ground only
if $A$ and $B$ are both 1, and two pMOS in parallel reach the supply if either
is 0. The NOR (Figure~\ref{fig:el:nor}) is the dual: parallel becomes series
and series becomes parallel, exactly as De Morgan's laws exchange AND and OR.
In every CMOS gate the network of pMOS that pulls the output up and the
network of nMOS that pulls it down are each other's dual, in the sense of the
duality principle of Chapter~\ref{ch:boole1}. Each of the two gates has four
transistors; the reader should check the four combinations of the inputs on
the drawings.

\subsection{What gates cost, in transistors}

\begin{table}[htb]
\centering\small
\caption{The cost of some circuits in conventional static CMOS; the counts
vary with the style of design.}\label{tab:el:cost}
\begin{tabular}{@{}lll@{}}
\toprule
\textbf{Circuit} & \textbf{Transistors} & \textbf{Remark}\\
\midrule
NOT & 2 & the CMOS inverter\\
NAND, NOR (2 inputs) & 4 & the natural gates of CMOS\\
AND, OR (2 inputs) & 6 & a NAND or a NOR followed by an inverter\\
XOR (2 inputs) & about 8 to 12 & depending on the design\\
full adder & about 28 & the classic static CMOS adder\\
SRAM cell, one bit & 6 & two inverters in a loop and two access transistors\\
D flip-flop & about 20 & two latches, master and slave\\
\bottomrule
\end{tabular}
\end{table}

An AND costs more than a NAND (Table~\ref{tab:el:cost}): in CMOS the
inverting gates come first. This is why NAND is the universal gate of
Chapter~\ref{ch:boole2} in practice as well as in theory. Minimising a
function means saving transistors, and therefore area, energy and time. The
adder behind the instruction \ins{add}, thirty-two full adders, is about 900
transistors.

\subsection{Switching takes time and energy}

The output of a gate drives the inputs of the next gates, which are small
capacitors. To change level the output must charge or discharge them through
the resistance of its transistors: this is the \term{propagation delay} of the
gate. And every change of level costs energy: the power spent in switching is
about
\[ P_{dyn} \approx C\cdot V^{2}\cdot f, \]
where $C$ is the capacitance switched, $V$ the supply voltage and $f$ the
frequency of switching. Lowering the supply voltage saves the most, because
it enters squared: it is the reason why chips moved from 5~V to 3.3~V and to
about 1~V inside. Raising the clock frequency raises the consumption in
proportion: speed is paid for in heat. The propagation delay returns in
Chapter~\ref{ch:combinational} as the critical path, and in
Chapter~\ref{ch:sequential} as the reason for the clock.

\subsection{Scales of integration}

\begin{table}[htb]
\centering\small
\caption{Scales of integration.}\label{tab:el:scales}
\begin{tabularx}{\textwidth}{@{}llL@{}}
\toprule
\textbf{Scale} & \textbf{Transistors per chip} & \textbf{Examples}\\
\midrule
SSI, small & up to about 100 & a few gates: the 74 series, such as four NANDs in one package\\
MSI, medium & hundreds & decoders, multiplexers, counters, registers: the machines of Chapters \ref{ch:combinational} to \ref{ch:registers}\\
LSI, large & thousands & the first microprocessors and memories: Intel 4004\\
VLSI, very large & hundreds of thousands and more & the processors of the 1980s and 1990s\\
today & tens of thousands of millions & systems on a chip: processor, memory, radio and peripherals together\\
\bottomrule
\end{tabularx}
\end{table}

The ESP32-C6 of the kit is a system on a chip in a package of a few
millimetres (Table~\ref{tab:el:scales}). The chips of the 74 series reappear
in the next chapters as the historical form of each building block.

% ----------------------------------------------------------------------------
\section{Back to the machine}

\subsection{Pull-up: an input that is never left floating}

An input connected to nothing is \term{floating}: it picks up noise and reads
0 or 1 at random. A \term{pull-up resistor} ties the pin to 3.3~V
(Figure~\ref{fig:el:pullup}). With the button released no current flows, and
the pin reads a clean 1. With the button pressed the switch joins the pin to
ground: 0~V, a clean 0, while a current of
$3.3\ \mathrm{V}/45\ \mathrm{k}\Omega = 73\ \mu\mathrm{A}$ flows through the
resistor, too small to matter. The circuit is the voltage divider of
(\ref{eq:el:div}) with a lower resistor that is either infinite or zero.

\begin{figure}[htb]
\centering
\begin{circuitikz}[scale=0.85,transform shape]
  \foreach \x/\sw/\v/\cap in {0/nos/{3.3 V: reads 1}/released, 6.5/ncs/{0 V: reads 0}/pressed}{
    \draw (\x,3.6) node[vcc]{3.3 V} to[R,l_=45<\kilo\ohm>] (\x,2) coordinate (p\sw);
    \draw (p\sw) to[\sw] (\x,0.2) node[ground]{};
    \draw (p\sw) node[circ]{} -- ++(1.2,0) node[buffer,anchor=in](b\sw){};
    \draw (b\sw.out) node[right]{\small \v};
    \node[lbl] at (\x+0.6,-0.9) {button \cap};
    \node[lbl,anchor=east] at (\x-0.1,2) {pin};
  }
\end{circuitikz}
\caption{The button of the kit with its pull-up resistor. The triangle is the
input buffer of the pin.}
\label{fig:el:pullup}
\end{figure}

The firmware enables the internal pull-up of every button pin, reads the pins
and inverts the bits, so that in the program pressed is 1. A signal whose
active state is the low level is said to use \term{active-low} logic.

\subsection{From the finger to the bit}

The questions of the warm-up now have answers, and Table~\ref{tab:el:finger}
puts them in order. Some of its steps belong to chapters still to come.

\begin{table}[htb]
\centering\small
\caption{From the finger to the bit.}\label{tab:el:finger}
\begin{tabularx}{\textwidth}{@{}cLP{0.34\textwidth}@{}}
\toprule
\textbf{Step} & \textbf{What happens} & \textbf{Where it is explained}\\
\midrule
1 & the finger closes the switch between GPIO18 and ground & this chapter: the circuit\\
2 & the pin falls from 3.3 V to 0 V, in well under a microsecond & this chapter: $R\cdot C$\\
3 & the input buffer sees a voltage below 0.825 V: a logic 0 & this chapter: logic levels\\
4 & the level is sampled into a flip-flop of the input register & Chapter \ref{ch:sequential}\\
5 & the firmware reads that register with a load instruction & Chapter \ref{ch:memory}: memory-mapped I/O\\
6 & it inverts the bit and packs it with the others into \reg{a0} & Chapters \ref{ch:boole1} and \ref{ch:io}\\
7 & the program tests it: \code{andi t3, a0, 0x01} & Chapter \ref{ch:firstcontact}\\
\bottomrule
\end{tabularx}
\end{table}

The contacts of a real switch \emph{bounce} for a few milliseconds when they
close. Reading the buttons once per frame, every few tens of milliseconds,
hides the bounce (Chapter~\ref{ch:io}).

\subsection{The kit as a circuit}

The kit of Figure~\ref{fig:setup:kit} can now be read as a circuit. A
\term{breadboard} has holes joined in rows underneath, so that components
connect without soldering. Two long buses carry the supply, 3.3~V and ground;
a separate bus carries 5~V for the audio amplifiers. Every part shares the
same ground: without a common reference the voltages of two parts mean nothing
to each other. The buttons go from a pin to ground; the display and the
amplifiers receive digital signals from other pins. An output pin can light
an LED through a resistor; an input pin only listens.

\begin{prgwarn}
A pin of the ESP32-C6 can supply some tens of milliamperes: an LED needs its
series resistor. Never join a pin that is driven high directly to ground.
\end{prgwarn}

% ----------------------------------------------------------------------------
\section{Check your understanding}

\begin{quiz}{An LED drops 2 V and must carry 5 mA from a 3.3 V pin. Which resistor is needed?}
\item 26 $\Omega$
\item 260 $\Omega$
\item 400 $\Omega$
\item 660 $\Omega$
\end{quiz}

\begin{quiz}{In a CMOS inverter with the input at 0:}
\item both transistors conduct
\item the pMOS conducts and the output is 1
\item the nMOS conducts and the output is 0
\item neither conducts and the output floats
\end{quiz}

\begin{quiz}{How many transistors does a two-input AND need in static CMOS?}
\item 2
\item 4
\item 6
\item 8
\end{quiz}

\begin{quiz}{Why does the button input of the kit need a pull-up resistor?}
\item To limit the current of the LED
\item To make the pin read a defined 1 when the button is released
\item To make the button faster
\item To protect the chip from 5 V
\end{quiz}

\section{Exercises}

\begin{exercise}[paper first, then the kit]
\begin{enumerate}[nosep]
\item A blue LED drops 3.0 V: which resistor gives 2 mA from 3.3 V? What power
  does the resistor dissipate?
\item A divider with $R_1=10$ k$\Omega$ and $R_2=20$ k$\Omega$ on 3.3 V:
  compute $V_{out}$. Is it read as 0, as 1, or is it undefined, by an input
  of the kit?
\item Draw a three-input NAND in CMOS and count its transistors.
\item Check the drawing of the NOR for the four combinations of the inputs,
  marking which transistors conduct.
\item On \prg, run the lamps cartridge and bridge a button pin to ground with a
  wire: explain what you see, and what would happen with the pull-up disabled.
\end{enumerate}
\expected{On paper: 150 $\Omega$ and 0.6 mW; $V_{out}=2.2$ V, in the
undefined band; six transistors, three pMOS in parallel and three nMOS in
series. On the kit: the lamp lights exactly as if the button were pressed.}
\end{exercise}

\section{Glossary}

\begin{glossarytable}
voltage & tensione\\
current & corrente\\
resistor / resistance & resistore / resistenza\\
capacitor & condensatore\\
ground & massa\\
semiconductor, doping & semiconduttore, drogaggio\\
gate, source, drain & gate, source, drain\\
logic level & livello logico\\
noise margin & margine di rumore\\
pull-up resistor & resistenza di pull-up\\
floating input & ingresso flottante\\
propagation delay & ritardo di propagazione\\
integrated circuit & circuito integrato\\
breadboard & basetta sperimentale (breadboard)\\
\end{glossarytable}

\section{Conclusions and reading}

Voltage pushes, current flows, resistance hinders: $V=R\cdot I$ describes
most of what a digital designer needs. A MOS transistor is a switch
controlled by a voltage, and nMOS and pMOS are complementary. A bit is a
range of voltages; noise margins and regeneration are what make digital
circuits reliable. CMOS gates are pairs of networks of switches, and the
inverting gates NAND and NOR are the cheapest. Switching costs time
($R\cdot C$) and energy ($C\cdot V^{2}\cdot f$): the physical limits behind
every clock.

\paragraph{Reading.} Charles Petzold, \emph{Code}, the chapters from the
flashlight to logic gates; \citet{patterson2020}, sections 1.5 and 1.7 and
appendix~A; \fname{docs/hardware/hardware.md} in the \prg repository.
